\documentclass[12pt,oneside]{book}

% ============================================================
% METADATOS
% ============================================================
\newcommand{\universidad}{Universidad de Buenos Aires (UBA)}
\newcommand{\facultad}{Facultad de Derecho}
\newcommand{\materia}{Derecho Civil}
\newcommand{\titulo}{El Estado Civil — Posición de la Persona en la Familia y la Sociedad}
\newcommand{\autor}{Isaac Edgar Camacho Ocampo}
\newcommand{\anio}{2026}

% ============================================================
% PAQUETES
% ============================================================
\usepackage[utf8]{inputenc}
\usepackage[T1]{fontenc}
\usepackage[spanish,es-nodecimaldot]{babel}

\usepackage[
    top=2.5cm,
    bottom=2.5cm,
    left=3cm,
    right=2.5cm,
    headheight=15pt
]{geometry}

\usepackage{amsmath}
\usepackage{amssymb}
\usepackage{mathtools}

\usepackage{graphicx}
\usepackage{float}
\usepackage{caption}
\usepackage{subcaption}

\usepackage{booktabs}
\usepackage{array}
\usepackage{longtable}
\usepackage{tabularx}

\usepackage{enumitem}

\usepackage{xcolor}
\usepackage[most]{tcolorbox}

\usepackage{fancyhdr}
\usepackage{ragged2e}

% ============================================================
% RUTAS DE IMÁGENES
% ============================================================
\graphicspath{
    {./img/}
    {./graficos/}
    {./figuras/}
}

% ============================================================
% CONFIGURACIÓN
% ============================================================
\setcounter{tocdepth}{2}

\setlength{\parindent}{0pt}
\setlength{\parskip}{0.5em}

\setlist[itemize]{
    topsep=0.3em,
    itemsep=0.2em
}

\setlist[enumerate]{
    topsep=0.3em,
    itemsep=0.2em
}

\clubpenalty=10000
\widowpenalty=10000

% ============================================================
% COMANDOS MATEMÁTICOS
% ============================================================
\newcommand{\abs}[1]{\left|#1\right|}
\newcommand{\norm}[1]{\left\lVert#1\right\rVert}
\newcommand{\vect}[1]{\mathbf{#1}}
\newcommand{\deriv}[2]{\frac{d#1}{d#2}}
\newcommand{\pderiv}[2]{\frac{\partial#1}{\partial#2}}

% ============================================================
% ENTORNOS / CAJAS ACADÉMICAS
% ============================================================
\newtcolorbox{definition}[1]{
    enhanced,
    breakable,
    title={Definición: #1},
    fonttitle=\bfseries,
    colback=blue!3,
    colframe=blue!50!black,
    boxrule=0.8pt,
    arc=2mm
}

\newtcolorbox{important}{
    enhanced,
    breakable,
    title={Importante},
    fonttitle=\bfseries,
    colback=yellow!8,
    colframe=orange!70!black,
    boxrule=0.8pt,
    arc=2mm
}

\newtcolorbox{example}[1]{
    enhanced,
    breakable,
    title={Ejemplo: #1},
    fonttitle=\bfseries,
    colback=green!4,
    colframe=green!40!black,
    boxrule=0.8pt,
    arc=2mm
}

\newtcolorbox{formula}[1]{
    enhanced,
    breakable,
    title={Fórmula / Regla: #1},
    fonttitle=\bfseries,
    colback=purple!4,
    colframe=purple!50!black,
    boxrule=0.8pt,
    arc=2mm
}

\newtcolorbox{exercise}[1]{
    enhanced,
    breakable,
    title={Ejercicio: #1},
    fonttitle=\bfseries,
    colback=gray!5,
    colframe=gray!60!black,
    boxrule=0.8pt,
    arc=2mm
}

\newtcolorbox{note}{
    enhanced,
    breakable,
    title={Nota},
    fonttitle=\bfseries,
    colback=white,
    colframe=gray!50,
    boxrule=0.6pt,
    arc=2mm
}

% ============================================================
% CONFIGURACIÓN FINAL: ENCABEZADOS Y PIES, HIPERVÍNCULOS
% ============================================================
\pagestyle{fancy}
\fancyhf{}

\fancyhead[L]{\nouppercase{\leftmark}}
\fancyhead[R]{\materia}

\fancyfoot[C]{\thepage}

\renewcommand{\headrulewidth}{0.4pt}
\renewcommand{\footrulewidth}{0pt}

\fancypagestyle{plain}{
    \fancyhf{}
    \fancyfoot[C]{\thepage}
    \renewcommand{\headrulewidth}{0pt}
}

\usepackage[
    colorlinks=true,
    linkcolor=blue!50!black,
    citecolor=green!40!black,
    urlcolor=blue!60!black,
    pdftitle={\titulo},
    pdfauthor={\autor},
    pdfsubject={Apuntes universitarios},
    bookmarks=true
]{hyperref}

% ============================================================
% DOCUMENTO
% ============================================================
\begin{document}

% ------------------------------------------------------------
% PORTADA
% ------------------------------------------------------------
\begin{titlepage}

\thispagestyle{empty}

\begin{center}

\vspace*{1cm}

{\large \textsc{\universidad}}

\vspace{0.2cm}

{\large \textsc{\facultad}}

\vspace{0.2cm}

{\large Introducción — Parte General}

\vspace{3cm}

{\Huge \textbf{\materia}}

\vspace{1cm}

{\LARGE \titulo}

\vspace{0.8cm}

\textit{Comprender los conceptos es el primer paso para convertir el estudio en conocimiento.}

\vspace{1.2cm}

\rule{0.70\textwidth}{0.5pt}

\vspace{0.5cm}

{\large Apuntes de estudio}

\vspace{0.5cm}

\rule{0.70\textwidth}{0.5pt}

\vfill

{\large \autor}

\vspace{0.3cm}

{\large Buenos Aires, Argentina}

\vspace{0.2cm}

{\large \anio}

\end{center}

\vfill

\begin{flushleft}
\textit{Este es un modesto aporte a los estudiantes de ``Derecho Civil'' no pretende sustituir el material oficial de la materia.}
\end{flushleft}

\end{titlepage}

% ------------------------------------------------------------
% ÍNDICE
% ------------------------------------------------------------
\tableofcontents

% ============================================================
% CAPÍTULO 1 — EL ESTADO CIVIL COMO ATRIBUTO
% ============================================================
\chapter{El estado civil como atributo de la persona}

Este capítulo desarrolla el estado civil como atributo de la persona humana: su concepto, sus antecedentes históricos en el derecho romano y su incidencia en otros institutos del derecho civil.

\section{Concepto de estado civil}

El término ``estado'' es polisémico. Habitualmente, cuando se habla de ``estado'' en derecho, se piensa en la organización política de una nación, es decir, en la organización política que lleva adelante la autoridad de un país. Sin embargo, esa es solo una de las acepciones del término.

\begin{definition}{Estado civil}
En derecho civil, el \textbf{estado} es la posición que la persona ocupa en la sociedad y en la familia. Se lo denomina habitualmente \textbf{estado civil}, noción que se asocia de inmediato con las categorías de casado, soltero o divorciado —más vinculadas a lo matrimonial—, pero que en realidad también comprende la \textbf{filiación} (la relación entre padres e hijos, la situación de paternidad y maternidad) y el \textbf{parentesco}.
\end{definition}

\section{Antecedentes históricos en el derecho romano}

En Roma, el estatus era determinante para fijar el plexo de derechos, es decir, el rango de derechos que tenía una persona. Existían tres estatus:

\begin{table}[H]
\centering
\begin{tabularx}{\textwidth}{l X}
\toprule
\textbf{Estatus} & \textbf{Contenido} \\
\midrule
\textit{Status libertatis} & Definía si la persona era libre o esclava. Si era libre, tenía una cantidad de derechos que no tenía el esclavo. \\
\textit{Status civitatis} & El estado en relación con la ciudad; definía si la persona era ciudadana romana o extranjera. Los ciudadanos romanos tenían derechos que no tenían los extranjeros. \\
\textit{Status familiae} & El estado en relación con la familia. La cabeza de la familia era el \textit{paterfamilias}, que tenía derechos que no tenían los restantes miembros de la familia. \\
\bottomrule
\end{tabularx}
\end{table}

El sistema vinculado al \textit{status libertatis} —es decir, la distinción entre personas libres y esclavas— ha sido abolido.

\section{El estado en la era de los derechos humanos}

Con el tiempo, la división en estatus que determinaba mayor o menor cantidad de derechos se fue desvaneciendo. En la era de los derechos humanos se entiende que no pueden hacerse distinciones de ese tipo entre las personas: todos los seres humanos son iguales en derechos y en condiciones. Por esta razón, el estado pierde relevancia en cuanto a la determinación de derechos.

\begin{important}
A lo sumo, reaparecen ciertas formas que podrían asimilarse a estados particulares, pero que apuntan, en realidad, a una mayor protección para los sectores más vulnerables. Es el caso de la Convención sobre los Derechos de las Personas con Discapacidad: cuando una persona presenta una discapacidad, tiene un conjunto mayor de protecciones y de derechos, porque se la considera vulnerable y se busca favorecer su plena inclusión y su pleno desarrollo en la vida social. Esto, sin embargo, no constituye el concepto de estado tradicional que se ha estudiado siempre: es una noción distinta.

La noción de estado civil se reduce, en definitiva, a la posición de la persona en la familia y en la sociedad.
\end{important}

\section{Importancia e incidencia del estado civil}

El estado civil tiene incidencia en otros institutos del derecho civil, lo que lo convierte en una herramienta conceptual relevante más allá de sí misma:

\begin{itemize}
    \item \textbf{Domicilio}: las personas incapaces, especialmente los hijos, mientras dura la responsabilidad parental, tienen su domicilio en el de sus padres.
    \item \textbf{Capacidad}: por ejemplo, los cónyuges no pueden contratar entre sí cuando están bajo el régimen de comunidad.
    \item \textbf{Apellido de los hijos}: los hijos matrimoniales tendrán el apellido de sus padres; si la filiación es extramatrimonial, tendrán el apellido solamente del progenitor que los reconoció, si tuvieran una sola filiación determinada.
    \item \textbf{Alimentos}: la determinación del estado incide en la obligación de prestar alimentos entre parientes, lo cual resulta especialmente relevante en el marco de la responsabilidad parental.
    \item \textbf{Otras obligaciones}, como la de informar ciertas situaciones cuando existe una relación de familia.
\end{itemize}

\begin{important}
Existen condiciones prohibidas en el Código vinculadas con el estado civil: no se puede imponer a alguien, como condición de un contrato o de una donación, que cambie o que no cambie su estado civil —por ejemplo, que no se case con determinada persona o que no se case nunca—, ni imponer como condición para donar una casa que la persona nunca contraiga matrimonio. Este tipo de condiciones resulta contrario a la libertad básica de las personas, razón por la cual está expresamente prohibido por el Código.
\end{important}

\begin{formula}{Artículo 344, Código Civil y Comercial}
Está prohibida la condición que afecte de modo grave la libertad de las acciones de la persona, como la de elegir un domicilio o una religión, o decidir sobre su estado civil —por ejemplo, no casarse, o casarse o no con determinada persona—.
\end{formula}

\section{El estado como atributo de la persona}

El estado civil interesa al derecho civil dentro de la parte general, en el grupo de los atributos de la persona: la persona humana tiene nombre, tiene una capacidad o aptitud para entablar relaciones jurídicas, tiene un patrimonio y también tiene un estado. Resulta importante conocer cuál es la ubicación de la persona en la familia y en la sociedad, precisamente por las razones expuestas: el estado tiene incidencia con los demás atributos y con deberes jurídicos propiamente civiles.

% ============================================================
% CAPÍTULO 2 — LAS ACCIONES DE ESTADO
% ============================================================
\chapter{Las acciones de estado}

Este capítulo desarrolla las herramientas procesales para discutir judicialmente el estado civil: las acciones de reclamación y de impugnación de la filiación.

\section{Concepto de acción de estado}

\begin{definition}{Acción de estado}
Las \textbf{acciones de estado} son acciones en sentido técnico jurídico: posibilidades de iniciar reclamos en sede judicial, orientados o bien a la \textbf{reclamación}, o bien a la \textbf{impugnación}, de un estado determinado. El caso típico en el que se desarrollan estas acciones es el de la \textbf{filiación}.
\end{definition}

\section{La acción de reclamación}

En la acción de reclamación, la persona busca ser emplazada en un estado determinado. Esto es lo típico que sucede cuando una persona no ha sido reconocida por uno de sus progenitores.

\begin{example}{Acción de reclamación de filiación}
Una persona que no ha sido reconocida por uno de sus progenitores puede iniciar una acción de reclamación: presenta las pruebas y ofrece un estudio genético para iniciar la acción. La parte demandada —la persona reclamada— contesta la demanda y puede no reconocer el vínculo. Se genera así un contradictorio en el que la parte reclamada ofrece pruebas, se las controvierte, se realiza el estudio genético, y finalmente se determina si la persona es o no hija mediante una sentencia judicial.
\end{example}

Este tipo de procesos se profundiza específicamente en la materia de derecho de familia; a los fines de esta introducción basta con comprender que existen acciones vinculadas con el estado de la persona, sin entrar en mayor detalle procesal.

\section{La acción de impugnación}

La acción de impugnación es la segunda gran categoría: se trata de una acción orientada a desplazar a una persona de un estado determinado.

\begin{example}{Acción de impugnación de la filiación}
Cuando una persona siente que quien pasa por su padre o por su madre no es, en realidad, su verdadero progenitor, puede iniciar una acción orientada a impugnar la filiación.
\end{example}

\begin{important}
A veces ocurre que una persona reconoce a un hijo creyendo que le está haciendo un bien, lo cual no es así en realidad. A este fenómeno se lo denomina \textbf{reconocimiento complaciente}. Es problemático porque no constituye la vía por la cual jurídicamente corresponde establecer la filiación: la filiación es muy importante para la identidad de la persona, pues hace al conocimiento de sus orígenes, de modo que no resulta irrelevante que alguien efectúe esa declaración sin que corresponda a la verdad biológica.
\end{important}

Cuando la persona reconocida no es verdaderamente el progenitor —siempre en el ámbito de la filiación por naturaleza—, se le impugna la paternidad o la maternidad: se sostiene que quien pasa por progenitor, en realidad, no lo es, y se inicia una acción de impugnación. Esta acción también puede ser iniciada por el propio interesado.

El Código regula distintas variantes sobre este punto, que exceden el alcance de esta introducción; basta con retener la existencia de estos dos tipos de acciones: de reclamación y de impugnación.

\section{El resultado de la prueba en las acciones de estado}

Como en todo proceso judicial, las acciones de estado pueden concluir exitosamente o no. Si el juez de primera instancia no hiciera lugar a las pruebas por considerar que lo alegado no está suficientemente probado, siempre existe la posibilidad de apelar a segunda instancia. En segunda instancia interviene generalmente un tribunal colegiado —al menos tres jueces, en la Capital Federal—, por lo que resulta difícil que en ninguna de las dos instancias se desconozcan las pruebas ofrecidas.

Existen, no obstante, pruebas más o menos fuertes: la prueba biológica y la prueba genética son especialmente importantes y decisivas, al menos en la filiación por naturaleza.

\begin{important}
Si se parte de un principio de verdad biológica, la ley presupone que el vínculo filiatorio se establece por el origen genético. Sin embargo, la filiación es un instituto complejo, compuesto también de elementos estáticos, como la posesión de estado, que se desarrolla en el capítulo siguiente.
\end{important}

% ============================================================
% CAPÍTULO 3 — LA POSESIÓN DE ESTADO
% ============================================================
\chapter{La posesión de estado}

Este capítulo estudia un concepto central y más sutil que el del estado formal: la posesión de estado, es decir, vivir un estado en los hechos, más allá del título, junto con sus cuatro aplicaciones concretas en el Código Civil y Comercial.

\section{Concepto de posesión de estado: la analogía con la posesión de las cosas}

La posesión de estado es un concepto analogable a la posesión de las cosas en el derecho civil. En el Código Civil y Comercial, la posesión se define como la situación de quien detenta una cosa con espíritu de ser el dueño de esa cosa, con independencia de si tiene o no el título correspondiente; a estas situaciones el Código las denomina relaciones de poder.

\begin{example}{La posesión de cosas y la prescripción adquisitiva}
La posesión tiene mucha importancia, por ejemplo, en los casos de prescripción adquisitiva: el poseedor que, durante un tiempo determinado —veinte años—, posee un inmueble como si fuera el dueño, puede iniciar una acción para ser tenido por dueño ante el Registro de la Propiedad Inmueble, precisamente porque poseyó el bien de manera ininterrumpida durante ese plazo. El Código denomina a esta figura \textbf{prescripción adquisitiva}.
\end{example}

\begin{definition}{Posesión de estado}
La \textbf{posesión de estado} es el goce de hecho de un estado determinado, con independencia de si la persona detenta o no el título correspondiente.
\end{definition}

\begin{important}
En la mayoría de los casos, título y posesión coinciden: si una persona tiene el título y, además, tiene la posesión —es decir, detenta en los hechos el estado—, no hay problema. Pero se trata de dos cosas distintas. Por ejemplo, en el estado de hijo: el \textbf{título} de hijo es estar anotado como hijo de determinada persona en el Registro Civil; pero, además, en los hechos, esa persona puede ser tratada como hijo por quienes figuran como sus padres, es decir, puede \textbf{poseer} el estado de hijo.
\end{important}

\section{Los elementos históricos de la posesión de estado: \textit{nomen}, \textit{tractatus} y \textit{fama}}

En Roma, la posesión de estado se configuraba a partir de tres elementos:

\begin{table}[H]
\centering
\begin{tabularx}{\textwidth}{l X}
\toprule
\textbf{Elemento} & \textbf{Contenido} \\
\midrule
\textit{Nomen} & El nombre, es decir, ser llamado como hijo (en el caso de la filiación). \\
\textit{Tractatus} & El trato, es decir, ser tratado como hijo. \\
\textit{Fama} & Ser conocido por los demás como hijo. \\
\bottomrule
\end{tabularx}
\end{table}

En la actualidad, el elemento más importante es el \textbf{trato} (\textit{tractatus}).

\begin{example}{Cómo se acredita el trato como hijo}
El trato propio de la relación de filiación se acredita, en la práctica, mediante pruebas del vínculo cotidiano: la fotografía de la celebración de un cumpleaños; el pago de la obra social; el pago de la cuota de un colegio, del club o de un profesor particular; o el hecho de llevar a la persona de vacaciones, con la fotografía correspondiente. Todas esas pruebas demuestran un trato de hijo.
\end{example}

La configuración de la posesión de un estado se produce cuando existe, en los hechos, el goce de esa situación de manera demostrable y probada, con independencia de si se tiene o no el título. Posesión de estado y título pueden o no coincidir: son dos cosas distintas. Una persona puede tener posesión de estado sin tener el título correspondiente, o viceversa; en la mayoría de los casos ambos coinciden, pero el Código atribuye especial relevancia a la posesión de estado en determinados supuestos, que se desarrollan a continuación.

\section{Las cuatro aplicaciones de la posesión de estado en el Código Civil y Comercial}

El Código Civil y Comercial vigente reconoce cuatro aplicaciones concretas de la noción de posesión de estado.

\subsection{Primera aplicación: la posesión de estado matrimonial}

Por regla general, la posesión de estado matrimonial no acredita el matrimonio, ya que el matrimonio es un acto jurídico solemne, que requiere el consentimiento expresado delante de un oficial del Registro Civil, con las formalidades y los impedimentos que señala el Código.

\begin{important}
Quienes simplemente conviven no pueden alegar que, por esa sola circunstancia, acceden a los beneficios del matrimonio. El Código Civil regula la unión convivencial, un tipo de vinculación familiar distinta, con su propia regulación específica; pero para acceder al estado de ``matrimonio'' se requiere haber celebrado el acto matrimonial, que es solemne.
\end{important}

La excepción central de esta primera aplicación se presenta cuando, al momento de celebrarse el matrimonio, existió algún problema de forma en el acta de celebración —por ejemplo, un error en el apellido de uno de los cónyuges—.

\begin{example}{Error formal en el acta de matrimonio}
Un cónyuge cuya pareja fallece se presenta al Registro Civil para iniciar la sucesión, y allí se le informa que existe un error en la confección del acta de matrimonio celebrado veinticinco años atrás, por lo cual el matrimonio sería nulo. Sin embargo, como transcurrieron veinticinco años y hubo posesión del estado de cónyuge durante todo ese tiempo, ese vicio formal queda saneado: las inobservancias de forma en la actuación no pueden hacerse valer, y el matrimonio se considera celebrado válidamente, a pesar del error en el acta.
\end{example}

\begin{formula}{Artículo 423, Código Civil y Comercial — Posesión de estado matrimonial}
La posesión de estado matrimonial no acredita el matrimonio. El matrimonio es un acto jurídico solemne. Sin embargo, si existió posesión de estado y se presenta un acta de celebración, los defectos formales del acta no pueden alegarse contra su existencia, si existe la posesión de estado.
\end{formula}

\subsection{Segunda aplicación: posesión de estado y reconocimiento de la filiación}

Esta es la aplicación más típica: la posesión de estado, debidamente acreditada en juicio, tiene el mismo valor que el reconocimiento, siempre que no sea desvirtuada por prueba en contrario sobre el nexo genético. Se aplica al caso de quien reclama de un progenitor la filiación mediante una acción de reclamación de filiación: si la persona prueba que el progenitor le pagaba la escuela, la trataba, convivía con ella y la llevaba de vacaciones, acredita la posesión de estado.

\begin{formula}{Artículo 584, Código Civil y Comercial — Posesión de estado y reconocimiento}
La posesión de estado debidamente acreditada en juicio tiene el mismo valor que el reconocimiento, siempre que no sea desvirtuada por prueba en contrario sobre el nexo genético.
\end{formula}

Esta regla vale como si el progenitor hubiera reconocido al hijo, pero queda desvirtuada si se acredita mediante prueba genética un nexo contrario; es decir, la prueba genética prevalece sobre la posesión de estado cuando ambas entran en conflicto.

\subsection{Tercera aplicación: el límite a la negación de la filiación presumida por la ley}

Cuando una persona está casada, se presume que los hijos nacidos durante el matrimonio son hijos de ambos cónyuges. El cónyuge que no dio a luz puede negar esa filiación dentro de un plazo determinado.

\begin{example}{Negación de la filiación presumida}
El supuesto más habitual es aquel en que, tras el nacimiento, el cónyuge que no dio a luz niega que el hijo sea suyo, alegando una infidelidad. La ley permite esta negación dentro de un plazo muy inmediato al nacimiento.
\end{example}

El límite central de esta tercera aplicación es que ese cónyuge no puede negar la filiación si hubo posesión de estado, es decir, si trató al hijo como propio. La posesión de estado impide, en ese caso, ejercer la negación de la filiación que en otro supuesto la ley autorizaría.

\begin{formula}{Artículo 591, Código Civil y Comercial — Límite a la negación de la filiación}
La ley autoriza al cónyuge a negar la filiación presumida apenas se produce el nacimiento, dentro de un plazo determinado. Pero no puede negarla si hubo posesión de estado, es decir, si trató al hijo como propio.
\end{formula}

\subsection{Cuarta aplicación: la adopción de personas mayores de edad}

En principio, siempre se adopta a una persona menor de edad; sin embargo, el Código permite adoptar también a una persona mayor de edad, de manera excepcional, cuando hubo posesión de estado de hijo, es decir, cuando esa persona fue tratada como hijo mientras era menor de edad.

\begin{important}
A la vuelta de los años, ya siendo la persona mayor de edad, puede adoptársela formalmente, porque fue tratada como hijo durante la minoridad. Esta figura puede tener un efecto sucesorio muy importante, porque el hijo adoptivo podría suceder a quien lo trató como su padre o madre durante el tiempo en que era menor de edad.
\end{important}

\begin{formula}{Artículo 597, Código Civil y Comercial — Adopción de personas mayores de edad}
Puede adoptarse a una persona mayor de edad, de manera excepcional, cuando hubo posesión de estado de hijo durante su minoridad, es decir, cuando fue tratada como hijo mientras era menor de edad.
\end{formula}

\subsection{Síntesis de las cuatro aplicaciones}

Las cuatro aplicaciones de la posesión de estado en el Código Civil y Comercial son, en síntesis: en relación con el matrimonio (artículo 423), en relación con el reconocimiento de la filiación (artículo 584), en relación con la posibilidad de negar la filiación presumida por la ley (artículo 591), y en relación con la adopción de personas mayores de edad (artículo 597).

% ============================================================
% CAPÍTULO 4 — LA PRUEBA DEL ESTADO CIVIL
% ============================================================
\chapter{La prueba del estado civil}

Este capítulo aborda cómo se prueba el estado civil ante el mundo jurídico, a través de las partidas del Registro del Estado Civil y Capacidad de las Personas y, cuando estas faltan, mediante otros medios.

\section{Las partidas como medio de prueba}

El nacimiento, sus circunstancias de tiempo y lugar, el sexo, el nombre y la filiación de las personas nacidas se prueban con las partidas del Registro Civil. Si una persona debe acreditar su filiación, tiene que acompañar la partida del Registro Civil correspondiente.

\begin{formula}{Regla general sobre la prueba del estado civil}
El nacimiento, sus circunstancias de tiempo y lugar, el sexo, el nombre y la filiación de las personas nacidas se prueban con las partidas del Registro Civil.
\end{formula}

\section{Organización del Registro Civil}

El Registro Civil se denomina técnicamente Registro del Estado Civil y Capacidad de las Personas. La forma de registración está regida por una ley nacional, pero la organización de cada registro depende de cada provincia, en virtud de la organización federal del país: cada provincia organiza su propio Registro Civil, respetando las disposiciones del Código y de la ley nacional.

\begin{formula}{Ley 26.413 — Registro del Estado Civil}
Regula la forma de registración de los hechos y actos relativos al estado civil de las personas en todo el país, uniformando la manera en que se registran tales hechos, aunque la organización de cada Registro Civil depende de cada provincia.
\end{formula}

El fundamento de esta exigencia de registración radica en que el estado civil está condicionado, está constituido a partir de hechos. El nacimiento es un hecho; el tiempo y el lugar en que ocurrió, así como el sexo al nacer, también lo son. El nombre, en cambio, se adjudica, y la filiación tampoco es un hecho puro, sino que queda determinada jurídicamente. No obstante, esos hechos o actos —como el matrimonio— se registran en el Registro del Estado Civil y Capacidad de las Personas, porque existe un interés público sustantivo en que la sociedad registre estos eventos.

Los grandes hechos y actos que el Registro Civil debe registrar son: la adopción, el matrimonio, el divorcio y, finalmente, el fallecimiento o defunción de la persona, que es el hecho que cierra el ciclo. De esta manera, el Registro Civil cuenta con partidas de nacimiento, partidas de matrimonio y partidas de defunción: estos son los tres grandes tipos de partida, en torno a los cuales se registran otros hechos con incidencia sobre ellos.

\section{El concepto de partida}

El concepto clave para entender la prueba del estado civil es el de partida.

\begin{definition}{Partida}
Las \textbf{partidas} son los asientos en el libro del Registro del Estado Civil y Capacidad de las Personas, en los que constan los hechos y actos constitutivos del estado: nacimiento, matrimonio y defunción.
\end{definition}

\begin{note}
A diferencia de lo que ocurre con las redes sociales —donde una persona puede consignar una edad que no es real—, la partida constituye la prueba formal del estado civil: es el documento donde consta que, en algún momento, ocurrió un hecho —que alguien nació, que un médico lo certificó, y que ese nacimiento se inscribió en el registro que lleva el Estado—.
\end{note}

\section{La partida como instrumento público}

\begin{important}
Las partidas son \textbf{instrumentos públicos}: tienen plena fe, tienen valor probatorio propio, y hacen plena fe de su contenido.
\end{important}

\section{Qué sucede cuando no hay partidas: la prueba supletoria}

La ausencia de partidas puede deberse a que no hay registro, a que falta la partida, a que falta el asiento, o a que ese asiento es nulo.

\begin{example}{El incendio de un Registro Civil}
Si el Registro Civil de una localidad sufriera un incendio que destruyera los libros correspondientes a determinado período, la persona cuyo nacimiento debía constar en esos libros puede recurrir a otros medios de prueba: puede acompañar otros certificados, como una partida vinculada a un hecho religioso —por ejemplo, haber sido bautizado— o una constancia de inscripción escolar en la que conste una determinada fecha de nacimiento. Mediante estas pruebas la persona puede acreditar la fecha de su nacimiento, o la fecha de defunción de otra persona, a los fines de la elaboración de la partida correspondiente.
\end{example}

\section{La diferencia entre la partida y el certificado médico}

\begin{note}
La partida no se confunde con el certificado médico de nacimiento o de defunción. El certificado médico es expedido por el médico respecto del hecho del nacimiento, y con él se concurre al Registro Civil, donde se confecciona la partida. La partida, en sí misma, es el documento que emana del Registro Civil; el certificado médico solamente le sirve a este para anotar el nacimiento. Por ello, lo que siempre debe ofrecerse como prueba es la partida emanada del Registro, y solo si esta no existiera corresponde recurrir a otros medios de prueba.
\end{note}

\section{El supuesto especial de la prueba supletoria del fallecimiento}

El artículo 98 contempla, en su segunda parte, un supuesto específico que, en rigor técnico, no es exactamente prueba supletoria del nacimiento o de la defunción de la manera ordinaria: el caso en que fallece una persona, pero no se encuentra el cadáver, o no puede identificarse, aunque no exista duda alguna de que esa persona falleció.

\begin{example}{El alpinista que cae en un lugar inaccesible}
Un alpinista sube una montaña peligrosa y cae por un precipicio, en un lugar absolutamente inaccesible, ante la vista de cinco compañeros que presencian la caída. Ninguno de ellos puede poner en duda que murió, porque la caída fue de cientos de metros de altura hacia un lugar inaccesible, y nunca se va a encontrar el cadáver. Este caso requiere un tratamiento especial: no puede declararse directamente fallecida a la persona, porque se necesita un médico que certifique la muerte, como en cualquier fallecimiento; tampoco corresponde acudir al procedimiento de presunción de fallecimiento, que es un proceso mucho más largo. Por ello, debe seguirse un procedimiento más breve, ante un juez, para disponer directamente la inscripción de la muerte.
\end{example}

\begin{example}{El accidente aéreo}
Otro supuesto del mismo tipo es el de un accidente aéreo en el que fallecen todos los ocupantes de la aeronave y es posible que los cadáveres nunca lleguen a identificarse individualmente. Si alguien prueba que una persona determinada se encontraba en ese vuelo, esa persona puede ser anotada como fallecida, y pueden iniciarse los trámites sucesorios correspondientes.
\end{example}

Este procedimiento no constituye un fallecimiento presunto, sino un fallecimiento indudable: las circunstancias de la desaparición deben ser tales que hagan tener la muerte por cierta. Una vez acreditado el fallecimiento por esta vía, se expide una partida de defunción común y corriente.

\begin{formula}{Artículo 98, Código Civil y Comercial — Prueba supletoria de fallecimiento sin hallazgo del cadáver}
Si el cadáver de una persona no es hallado, o no puede ser identificado, el juez puede tener por comprobada la muerte y disponer la pertinente inscripción en el registro, siempre que la desaparición se hubiera producido en circunstancias tales que la muerte deba ser tenida como cierta.
\end{formula}

% ============================================================
% RESUMEN ESTILO CORNELL
% ============================================================
\chapter*{Resumen estilo Cornell}
\addcontentsline{toc}{chapter}{Resumen estilo Cornell}

La siguiente tabla organiza los conceptos centrales de la clase en formato Cornell: a la izquierda, las preguntas o conceptos clave; a la derecha, su desarrollo o respuesta. Al final se incluye una síntesis integradora de todo el contenido.

\begin{longtable}{
    >{\raggedright\arraybackslash}p{0.28\textwidth}
    >{\raggedright\arraybackslash}p{0.62\textwidth}
}
\toprule
\textbf{Preguntas / palabras clave} & \textbf{Notas / respuestas} \\
\midrule
\endfirsthead

\toprule
\textbf{Preguntas / palabras clave} & \textbf{Notas / respuestas} \\
\midrule
\endhead

\bottomrule
\endfoot

\bottomrule
\endlastfoot

\textbf{¿Qué es el estado civil?} &
La posición que la persona ocupa en la sociedad y en la familia. Se asocia habitualmente con ser casado, soltero o divorciado, pero también comprende la filiación (paternidad, maternidad) y el parentesco. \\

\textbf{¿Qué eran los tres estatus del derecho romano?} &
El \textit{status libertatis} (libre o esclavo), el \textit{status civitatis} (ciudadano romano o extranjero) y el \textit{status familiae} (posición dentro de la familia, encabezada por el \textit{paterfamilias}). Determinaban el rango de derechos de cada persona. \\

\textbf{¿Por qué el estado perdió relevancia como fuente de desigualdad de derechos?} &
Porque en la era de los derechos humanos se entiende que todos los seres humanos son iguales en derechos y condiciones. Solo reaparecen protecciones especiales para grupos vulnerables, como las personas con discapacidad, sin que ello constituya el concepto clásico de estado. \\

\textbf{¿En qué institutos del derecho civil incide el estado?} &
En el domicilio (los incapaces tienen el de su representante), en la capacidad (por ejemplo, restricciones para contratar entre cónyuges), en el apellido de los hijos, y en la obligación alimentaria entre parientes. \\

\textbf{¿Qué condiciones vinculadas al estado civil están prohibidas?} &
No se puede imponer a alguien, como condición de un contrato o de una donación, que cambie o no cambie su estado civil (por ejemplo, que no se case nunca), porque afecta la libertad básica de las personas (artículo 344). \\

\textbf{¿Qué es una acción de estado?} &
Una acción judicial orientada a reclamar el emplazamiento en un estado determinado, o a impugnar (desplazar) a alguien de un estado que ostenta. El caso típico es la filiación. \\

\textbf{¿Qué busca la acción de reclamación?} &
Que la persona sea emplazada en un estado que no tiene reconocido; por ejemplo, que un hijo no reconocido por su progenitor obtenga, mediante prueba genética, el reconocimiento judicial de la filiación. \\

\textbf{¿Qué busca la acción de impugnación?} &
Desplazar a alguien de un estado que ostenta pero que no le corresponde realmente; por ejemplo, impugnar la paternidad de quien pasa por progenitor sin serlo biológicamente. \\

\textbf{¿Qué es el reconocimiento complaciente?} &
El reconocimiento de un hijo realizado creyendo erróneamente que se le hace un bien, sin que corresponda a la verdad biológica. No es la vía jurídicamente correcta para establecer la filiación, porque afecta el derecho a la identidad. \\

\textbf{¿Qué es la posesión de estado?} &
El goce de hecho de un estado determinado, con independencia de si la persona tiene o no el título correspondiente. Es un concepto análogo a la posesión de las cosas en derecho civil. \\

\textbf{¿Cuáles eran los tres elementos romanos de la posesión de estado?} &
\textit{Nomen} (el nombre, ser llamado hijo), \textit{tractatus} (el trato, ser tratado como hijo) y \textit{fama} (ser conocido por los demás como hijo). Hoy el elemento más importante es el trato. \\

\textbf{¿Qué efecto tiene la posesión de estado matrimonial ante un error formal en el acta?} &
Si hubo posesión de estado de cónyuges durante un tiempo prolongado, los defectos de forma del acta de matrimonio no pueden alegarse contra su existencia (artículo 423). \\

\textbf{¿Qué valor tiene la posesión de estado en un juicio de filiación?} &
Tiene el mismo valor que el reconocimiento, si está debidamente acreditada en juicio, salvo que sea desvirtuada por prueba genética en contrario (artículo 584). \\

\textbf{¿Puede un cónyuge negar la filiación de un hijo si lo trató como propio?} &
No. Aunque la ley permite negar la filiación presumida dentro de un plazo inmediato al nacimiento, no puede hacerlo si hubo posesión de estado, es decir, si trató al hijo como tal (artículo 591). \\

\textbf{¿Se puede adoptar a una persona mayor de edad?} &
Excepcionalmente sí, cuando hubo posesión de estado de hijo durante su minoridad. Puede tener efectos sucesorios importantes, ya que el adoptado podría suceder a quien lo trató como padre o madre (artículo 597). \\

\textbf{¿Cómo se prueba el estado civil?} &
A través de las partidas del Registro del Estado Civil y Capacidad de las Personas: partidas de nacimiento, de matrimonio y de defunción. \\

\textbf{¿Cómo se organiza el Registro Civil en la Argentina?} &
La forma de registración está regida por la Ley 26.413 a nivel nacional, pero cada provincia organiza su propio Registro Civil, en virtud de la organización federal del país. \\

\textbf{¿Qué son las partidas?} &
Los asientos en el libro del Registro del Estado Civil, donde constan los hechos y actos constitutivos del estado: nacimiento, matrimonio y defunción. Son instrumentos públicos: hacen plena fe de su contenido. \\

\textbf{¿La partida es lo mismo que el certificado médico de nacimiento?} &
No. El certificado médico es el documento que expide el médico y que sirve al Registro Civil para confeccionar la partida; la partida en sí es el documento que emana del propio Registro Civil. \\

\textbf{¿Qué se puede hacer si no existe la partida (por ejemplo, por un incendio en el Registro)?} &
Se puede recurrir a la prueba supletoria: otros certificados o constancias, como partidas religiosas (por ejemplo, de bautismo) o constancias escolares que registren la fecha de nacimiento. \\

\textbf{¿Cómo se inscribe la muerte de una persona cuando no se encuentra el cadáver?} &
Mediante un procedimiento judicial específico (artículo 98), distinto de la presunción de fallecimiento, cuando las circunstancias de la desaparición hacen tener la muerte por cierta (por ejemplo, una caída en un lugar inaccesible o un accidente aéreo). \\

\midrule

\multicolumn{2}{p{\dimexpr\textwidth-2\tabcolsep}}{
\textbf{Resumen:} El estado civil es la posición que la persona ocupa en la familia y en la sociedad; sus antecedentes romanos (\textit{status libertatis}, \textit{civitatis} y \textit{familiae}) perdieron relevancia como fuente de desigualdad de derechos en la era de los derechos humanos, aunque subsisten protecciones especiales para grupos vulnerables. El estado incide en el domicilio, la capacidad, el apellido y los alimentos, y no puede imponerse como condición de un acto jurídico (artículo 344). Las acciones de estado —reclamación e impugnación— permiten discutir judicialmente la filiación. La posesión de estado, goce de hecho de un estado con independencia del título, configurada históricamente por el nombre, el trato y la fama, tiene cuatro aplicaciones en el Código: sanea defectos formales del acta matrimonial (artículo 423), equivale al reconocimiento en juicio de filiación salvo prueba genética en contrario (artículo 584), impide negar la filiación presumida cuando hubo trato de hijo (artículo 591), y habilita la adopción excepcional de personas mayores de edad (artículo 597). Finalmente, el estado civil se prueba mediante las partidas del Registro del Estado Civil y Capacidad de las Personas —instrumentos públicos con plena fe—, distintas del certificado médico, y ante su ausencia el ordenamiento admite la prueba supletoria, incluido el supuesto especial del fallecimiento sin hallazgo del cadáver (artículo 98).
} \\

\end{longtable}

\end{document}
